```latex
\documentclass[11pt]{article}

\usepackage{amsmath,amssymb,amsthm,mathtools}
\usepackage[colorlinks=true,linkcolor=blue,citecolor=blue,urlcolor=blue]{hyperref}

\newtheorem{theorem}{Theorem}[section]
\newtheorem{lemma}[theorem]{Lemma}
\newtheorem{corollary}[theorem]{Corollary}
\theoremstyle{definition}
\newtheorem{definition}[theorem]{Definition}
\theoremstyle{remark}
\newtheorem{remark}[theorem]{Remark}

\newcommand{\Sthree}{S^{3}}
\newcommand{\HFhat}{\widehat{HF}}
\newcommand{\SpinC}{\mathrm{Spin}^{c}}
\newcommand{\ZZ}{\mathbb Z}

\title{Reducible surgeries on \texorpdfstring{$L$}{L}-space knots have at
most two prime summands}
\author{}
\date{}

\begin{document}

\maketitle

\begin{abstract}
We prove that every manifold obtained by Dehn surgery on an $L$-space knot
in $S^{3}$ has at most two prime summands.  In particular, a reducible
surgery on an $L$-space knot has exactly two prime summands.  The proof
combines the general three-summands theorem for Dehn surgery, the
Matignon--Sayari slope bound, the Hom--Lidman--Zufelt lower bound for
reducing slopes on $L$-space knots, and a correction-term obstruction at the
critical slope $2g(K)-1$.
\end{abstract}

\section{Statement and conventions}

For a knot $K\subset \Sthree$ and a slope $r\in \mathbb Q\cup\{\infty\}$,
write
\[
K(r)=\Sthree_r(K)
\]
for $r$-framed Dehn surgery on $K$.  All $3$-manifolds are closed and
oriented.  We use the convention that a \emph{positive $L$-space knot} is a
knot admitting a positive-slope surgery which is a Heegaard Floer $L$-space.
If a knot admits a negative $L$-space surgery, its mirror is a positive
$L$-space knot.  Since
\[
\Sthree_{-r}(\overline K)\cong -\Sthree_r(K),
\]
this operation preserves the number of prime summands.  Consequently, there
is no loss in proving the result for positive $L$-space knots.

The unknot is immediate: every surgery on it is $S^3$, $S^1\times S^2$, or a
lens space, and hence is prime.  Throughout the proof, $K$ is therefore a
nontrivial positive $L$-space knot of genus $g=g(K)$.

\begin{theorem}\label{thm:main}
Let $K\subset \Sthree$ be an $L$-space knot and let $r$ be any Dehn-surgery
slope.  Then $K(r)$ has at most two prime summands.  Equivalently, if
$K(r)$ is reducible, then it has exactly two prime summands.
\end{theorem}

This theorem does not assert the full cabling conjecture for $L$-space
knots: it controls only the number of prime summands, not the isotopy type
of $K$.

\section{Topological reduction theorems}

We first record the classical restrictions on reducible Dehn surgeries.

\begin{lemma}[Gordon--Luecke integrality theorem]
\label{lem:GL}
If Dehn surgery on a nontrivial knot in $\Sthree$ is reducible, then the
surgery slope is an integer.  Moreover, after mirroring if necessary, the
surgery may be assumed to have positive integral slope, and the surgered
manifold contains a lens-space summand.
\end{lemma}

\begin{proof}
These are the principal reducible-surgery results of Gordon and Luecke
\cite{GLReducible}.
\end{proof}

The next theorem is the general three-summands reduction theorem for Dehn
surgery.

\begin{lemma}[Howie--Sayari--Valdez S\'anchez]
\label{lem:three-summands}
Let $K$ be a knot in $\Sthree$ for which $K(r)$ is reducible.  Then $K(r)$
has at most three prime summands.  If it has exactly three prime summands,
then, after orienting so that $r=n>0$,
\[
K(n)\cong L(p,a)\# L(q,b)\# Y,
\]
where $p,q>1$ are relatively prime, $Y$ is a prime integer homology sphere,
and
\[
n=pq.
\]
\end{lemma}

\begin{proof}
This is the combined three-summands theorem of Howie, Sayari, and Valdez
S\'anchez.  A precise statement and discussion of the original papers
appears in the introduction to Howie's article
\emph{Can Dehn surgery yield three connected summands?}
\cite{HowieThree}; an accessible HTML rendering is at
\href{https://ar5iv.labs.arxiv.org/html/0908.2697}{ar5iv.labs.arxiv.org}.
The equality $n=pq$ follows by comparing first homology:
\[
\ZZ/n\cong H_1(K(n);\ZZ)
       \cong \ZZ/p\oplus \ZZ/q.
\]
\end{proof}

We also need the standard description of surgery at the cabling slope.

\begin{lemma}[Cabling-slope decomposition]
\label{lem:cable-decomposition}
Let $K=C_{p,q}(J)$ be a nontrivial $(p,q)$-cable of a knot $J\subset
\Sthree$, where $p>1$ is the winding number.  At the cabling slope $pq$,
\[
K(pq)\cong J(q/p)\#L(p,q),
\]
up to the conventional relabeling of the lens-space parameter.  In
particular, $K(pq)$ has at most two prime summands.
\end{lemma}

\begin{proof}
The cabling annulus splits the exterior of $K$ into the companion exterior
and a solid-torus piece.  Filling along the cabling slope produces the
displayed connected-sum decomposition; see Gordon's treatment of satellite
surgeries \cite{GordonSatellite}.

For the last assertion, the lens-space factor is prime.  If $J$ is the
unknot, then $J(q/p)$ is also a lens space, hence prime.  If $J$ is
nontrivial and $J(q/p)$ were reducible, Lemma~\ref{lem:GL} would imply that
$q/p$ is integral, contradicting $\gcd(p,q)=1$ and $p>1$.  Thus $J(q/p)$ is
irreducible.  Since it has finite first homology, it is not $S^1\times S^2$,
and hence it is prime.  Therefore the displayed manifold has at most two
prime summands.
\end{proof}

\section{The slope bounds}

The following theorem of Matignon and Sayari is the topological upper bound
for a non-cabling reducing slope.

\begin{lemma}[Matignon--Sayari]
\label{lem:MS}
Let $K$ be a nontrivial knot of genus $g$, and suppose that $K(n)$ is
reducible for an integer $n>0$.  Then either
\begin{enumerate}
\item $K$ is a cable knot and $n$ is its cabling slope, or
\item
\[
n\le 2g(K)-1.
\]
\end{enumerate}
\end{lemma}

\begin{proof}
This is Theorem~1.1 of Matignon--Sayari \cite{MatignonSayari}.  The
statement is also quoted and placed in the context of the cabling
conjecture in \cite{LSGB}; see the publication record at
\href{https://exa.ai/library/publication/zl1m8jt741b}{exa.ai}.
\end{proof}

For $L$-space knots there is a Heegaard Floer lower bound in the opposite
direction.

\begin{lemma}[Hom--Lidman--Zufelt]
\label{lem:HLZ}
Let $K$ be a positive $L$-space knot of genus $g$.  If $K(r)$ is reducible,
then $r$ is an integer and
\[
r\ge 2g-1.
\]
\end{lemma}

\begin{proof}
Integrality is Lemma~\ref{lem:GL}.  The lower bound is the
Hom--Lidman--Zufelt reducible-surgery theorem for $L$-space knots
\cite{HLZ15}.  Their argument uses the knot Floer surgery formula together
with the Gordon--Luecke lens-space summand.  The result and its role in the
cabling conjecture are also discussed in the recent treatment
\cite{RecentObstructions}, rendered at
\href{https://ar5iv.labs.arxiv.org/html/2209.01672}{ar5iv.labs.arxiv.org}.
\end{proof}

Combining the two bounds gives the critical slope.

\begin{corollary}\label{cor:critical-slope}
Suppose that $K$ is a positive $L$-space knot of genus $g$ and that $K(n)$
has three prime summands.  Then
\[
n=2g-1.
\]
\end{corollary}

\begin{proof}
By Lemma~\ref{lem:HLZ}, $n\ge 2g-1$.  If $K$ were a cable knot and $n$ its
cabling slope, Lemma~\ref{lem:cable-decomposition} would imply that $K(n)$
has at most two prime summands, contrary to hypothesis.  Hence the second
alternative of Lemma~\ref{lem:MS} holds, giving $n\le 2g-1$.
\end{proof}

\section{Correction-term preliminaries}

Let $K$ be a positive $L$-space knot of genus $g$.  For $s\ge 0$, let
$V_s(K)$ denote the Ozsv\'ath--Szab\'o--Rasmussen integer surgery
invariants.  For an $L$-space knot these equal the torsion coefficients of
the Alexander polynomial.  The standard identities we need are
\begin{equation}\label{eq:V-identities}
V_0(K)=g,\qquad V_g(K)=0,\qquad V_{g-1}(K)=1.
\end{equation}
Indeed, the normalized Alexander polynomial of an $L$-space knot has
rank-one knot Floer groups in every supporting Alexander grading, and its
top Alexander coefficient is $1$; the displayed identities are then the
standard torsion-coefficient formulas \cite{OSLens,NiWu}.

For a positive integer $n$, label the $\SpinC$ structures on $K(n)$ by
$i\in\ZZ/n$.  The absolutely graded surgery formula gives
\begin{equation}\label{eq:surgery-d}
d\bigl(K(n),\mathfrak s_i\bigr)
=
d\bigl(L(n,1),i\bigr)
-
2\max\{V_i(K),V_{n-i}(K)\},
\qquad 0\le i<n,
\end{equation}
where
\begin{equation}\label{eq:lens-d}
d\bigl(L(n,1),i\bigr)
=
\frac{(2i-n)^2-n}{4n}.
\end{equation}
See Ozsv\'ath--Szab\'o and Rasmussen \cite{OSAbs,OSLens}, and Ni--Wu
\cite{NiWu}.

We need only the following elementary universal bound for arbitrary lens
spaces.

\begin{lemma}\label{lem:lens-bound}
For every lens space $L(m,a)$ and every $\mathfrak s\in\SpinC(L(m,a))$,
\[
\left|d(L(m,a),\mathfrak s)\right|\le \frac m4.
\]
\end{lemma}

\begin{proof}
This is the standard correction-term estimate for lens spaces, obtained
from the continued-fraction plumbing formula for lens-space $d$-invariants
\cite{OSAbs}.
\end{proof}

Finally, correction terms are additive under connected sum:
\begin{equation}\label{eq:connected-sum-d}
d(Y_1\#Y_2,\mathfrak s_1\#\mathfrak s_2)
=
d(Y_1,\mathfrak s_1)+d(Y_2,\mathfrak s_2).
\end{equation}
See \cite{OSAbs}.

\section{The three-summand obstruction}

The following lemma is the new numerical step in the argument.

\begin{lemma}\label{lem:d-obstruction}
Let $K$ be a positive $L$-space knot of genus $g$, and put $n=2g-1$.
There do not exist integers $p,q>1$, a homology sphere $Y$, and lens-space
parameters $a,b$ such that
\[
K(n)\cong L(p,a)\#L(q,b)\#Y
\qquad\text{and}\qquad
n=pq.
\]
\end{lemma}

\begin{proof}
Suppose such a decomposition exists.  Since $n=2g-1$ is odd and $n=pq$,
both $p$ and $q$ are odd and hence
\begin{equation}\label{eq:pq-lower}
p,q\ge 3.
\end{equation}

Let $\mathfrak s_Y$ be the unique $\SpinC$ structure on the homology sphere
$Y$, and write
\[
d_Y=d(Y,\mathfrak s_Y).
\]
Every $\SpinC$ structure on $K(n)$ is a connected sum
\[
\mathfrak s=\mathfrak s_p\#\mathfrak s_q\#\mathfrak s_Y,
\]
and \eqref{eq:connected-sum-d} gives
\[
d(K(n),\mathfrak s)
=
d(L(p,a),\mathfrak s_p)
+
d(L(q,b),\mathfrak s_q)
+
d_Y.
\]

First consider the surgery label $i=0$.  By
\eqref{eq:V-identities}--\eqref{eq:lens-d},
\[
\begin{aligned}
d(K(n),\mathfrak s_0)
&=
\frac{n-1}{4}-2V_0(K)\\
&=
\frac{n-1}{4}-2g.
\end{aligned}
\]
Since $n=2g-1$, this becomes
\begin{equation}\label{eq:d0}
d(K(n),\mathfrak s_0)
=
\frac{n-1}{4}-(n+1)
=
-\frac{3n+5}{4}.
\end{equation}
Under the connected-sum identification, $\mathfrak s_0$ corresponds to some
pair $(\mathfrak s_p^0,\mathfrak s_q^0)$.  Thus
\[
d_Y
=
d(K(n),\mathfrak s_0)
-
d(L(p,a),\mathfrak s_p^0)
-
d(L(q,b),\mathfrak s_q^0).
\]
Lemma~\ref{lem:lens-bound} and \eqref{eq:d0} yield the upper bound
\begin{equation}\label{eq:dY-upper}
d_Y
\le
-\frac{3n+5}{4}+\frac p4+\frac q4
=
\frac{p+q-3n-5}{4}.
\end{equation}

Next consider the surgery label $i=g$.  Since $n-g=g-1$,
\[
\max\{V_g(K),V_{n-g}(K)\}
=
\max\{0,V_{g-1}(K)\}
=
1
\]
by \eqref{eq:V-identities}.  Moreover, because $2g-n=1$,
\[
d(L(n,1),g)
=
\frac{1-n}{4n}
>
-\frac14.
\]
Consequently \eqref{eq:surgery-d} gives
\begin{equation}\label{eq:dg-lower}
d(K(n),\mathfrak s_g)
>
-\frac14-2
=
-\frac94.
\end{equation}
If $\mathfrak s_g$ corresponds to
$(\mathfrak s_p^g,\mathfrak s_q^g)$, then
\[
d_Y
=
d(K(n),\mathfrak s_g)
-
d(L(p,a),\mathfrak s_p^g)
-
d(L(q,b),\mathfrak s_q^g).
\]
Using \eqref{eq:dg-lower} and Lemma~\ref{lem:lens-bound}, we obtain
\begin{equation}\label{eq:dY-lower}
d_Y
>
-\frac94-\frac p4-\frac q4
=
-\frac{p+q+9}{4}.
\end{equation}

Comparing \eqref{eq:dY-lower} and \eqref{eq:dY-upper} would imply
\[
-\frac{p+q+9}{4}
<
d_Y
\le
\frac{p+q-3n-5}{4},
\]
and therefore
\[
3n-4<2p+2q.
\]
But $n=pq$ and $p,q\ge3$, so
\[
\begin{aligned}
3n-4-2p-2q
&=
3pq-4-2p-2q\\
&=
p(3q-2)-2q-4\\
&\ge 3(3q-2)-2q-4\\
&=7q-10\ge 11>0.
\end{aligned}
\]
This contradiction proves the lemma.
\end{proof}

\section{Proof of the theorem}

\begin{proof}[Proof of Theorem~\ref{thm:main}]
If $K$ is the unknot, every surgery is prime.  Hence assume that $K$ is a
nontrivial positive $L$-space knot.

If $K(r)$ is irreducible, it has one prime summand.  Suppose instead that
$K(r)$ is reducible.  By Lemma~\ref{lem:GL} and Lemma~\ref{lem:HLZ}, $r$
is a positive integer, say $r=n$, and
\[
n\ge 2g-1.
\]
By Lemma~\ref{lem:three-summands}, $K(n)$ has at most three prime summands.
Assume for contradiction that it has exactly three.  Then
\[
K(n)\cong L(p,a)\#L(q,b)\#Y,
\]
where $p,q>1$ are relatively prime, $Y$ is a prime integer homology sphere,
and $n=pq$.

By Corollary~\ref{cor:critical-slope},
\[
n=2g-1.
\]
Thus all the hypotheses of Lemma~\ref{lem:d-obstruction} are satisfied.
That lemma says that no such connected-sum decomposition exists.  This is a
contradiction.  Therefore $K(n)$ does not have three prime summands, and
hence has at most two.

Since a reducible manifold has at least two prime summands, a reducible
surgery on an $L$-space knot has exactly two.
\end{proof}

\begin{remark}
The correction-term argument rules out a three-summand decomposition at the
critical slope $2g-1$.  The argument above does not prove that a
two-summand surgery is necessarily the surgery predicted by the full
cabling conjecture.
\end{remark}

\begin{thebibliography}{99}

\bibitem{GLReducible}
C.~McA. Gordon and J.~Luecke,
\emph{Reducible manifolds and Dehn surgery},
Topology \textbf{35} (1996), 385--409.

\bibitem{GordonSatellite}
C.~McA. Gordon,
\emph{Dehn surgery and satellite knots},
Transactions of the American Mathematical Society \textbf{275} (1983),
687--708.

\bibitem{HowieThree}
J.~Howie,
\emph{Can Dehn surgery yield three connected summands?},
arXiv:0908.2697.  Accessible rendering:
\href{https://ar5iv.labs.arxiv.org/html/0908.2697}{ar5iv.labs.arxiv.org}.

\bibitem{MatignonSayari}
D.~Matignon and N.~Sayari,
\emph{Longitudinal slope and Dehn fillings},
Hiroshima Mathematical Journal \textbf{33} (2003), 127--136.

\bibitem{LSGB}
\emph{L-space surgeries, genus bounds, and the cabling conjecture}.
Publication record and quoted Matignon--Sayari theorem:
\href{https://exa.ai/library/publication/zl1m8jt741b}{exa.ai}.

\bibitem{HLZ15}
J.~Hom, T.~Lidman, and N.~Zufelt,
\emph{Reducible surgeries and Heegaard Floer homology},
Mathematical Research Letters \textbf{22} (2015).

\bibitem{RecentObstructions}
\emph{Obstructing Reducible Surgeries: Slice Genus and Thickness Bounds},
arXiv:2209.01672.  Accessible rendering:
\href{https://ar5iv.labs.arxiv.org/html/2209.01672}{ar5iv.labs.arxiv.org}.

\bibitem{OSAbs}
P.~Ozsv\'ath and Z.~Szab\'o,
\emph{Absolutely graded Floer homologies and intersection forms for
four-manifolds with boundary},
Advances in Mathematics \textbf{173} (2003), 179--261.

\bibitem{OSLens}
P.~Ozsv\'ath and Z.~Szab\'o,
\emph{On knot Floer homology and lens space surgeries},
Topology \textbf{44} (2005), 1281--1300.

\bibitem{NiWu}
Y.~Ni and Z.~Wu,
\emph{Cosmetic surgeries on knots in $S^3$},
Journal f\"ur die reine und angewandte Mathematik \textbf{706} (2015),
1--17.

\end{thebibliography}

\end{document}
```